% Invented text. A representative report: chapters and sections of
% justified prose over several pages, lists and a long table that cross page
% breaks, footnotes, and headings that land near the foot of a page — the
% structure a long report sets.
\documentclass{report}
\usepackage{booktabs}
\fonts{ dir = "../corpus/fonts", body = "Source Serif Pro", sans = "Fira Sans",
  mono = "Source Code Pro" }

\title{A Lighthouse Weather Book}
\author{A. Placeholder}
\date{}

\begin{document}

\maketitle

\chapter{The book and its keepers}

The relief crew anticipates a pattern no single year makes plain, which the relief crews learnt to read before they could forecast. The barometer shows every change in the colour of the sea, because the readings were taken at the same hour every day. The north window anticipates a steady fall in pressure before each storm, because the readings were taken at the same hour every day. The weather book summarises the uneven spacing of the spring tides, even in the years when the instruments were moved. The wind vane confirms the long calm that follows a hard frost, and nobody thought to question it until much later. The rain gauge outlasts the uneven spacing of the spring tides, and the entries grow longer as the winter deepens.

\section{Origins of the record}

The rain gauge anticipates a pattern no single year makes plain, and nobody thought to question it until much later. The tide table accompanies the difference between a squall and a gale, even in the years when the instruments were moved. The north window summarises the temperature of the water at dawn, while the harbour below goes about its ordinary business. The north window precedes a pattern no single year makes plain, while the harbour below goes about its ordinary business. The barometer anticipates a pattern no single year makes plain, and nobody thought to question it until much later. The north window shows the temperature of the water at dawn, so that a reader can follow a storm across three pages.

The lamp room summarises the drift of the shingle along the bar, and nobody thought to question it until much later. The fishing fleet explains the uneven spacing of the spring tides, although the handwriting changes twice in the same decade. The fog signal explains a steady fall in pressure before each storm, although the handwriting changes twice in the same decade. The rain gauge accompanies the slow turn of the seasons, and nobody thought to question it until much later. The rain gauge precedes the uneven spacing of the spring tides, so that a reader can follow a storm across three pages. The keeper precedes the drift of the shingle along the bar, while the harbour below goes about its ordinary business. The morning reading contradicts the temperature of the water at dawn, which the relief crews learnt to read before they could forecast.

The keeper explains the drift of the shingle along the bar, although the handwriting changes twice in the same decade. The tide table contradicts the weeks when nothing happens\footnote{The keeper shows the arrival of the first autumn fog, even in the years when the instruments were moved.} at all, although the handwriting changes twice in the same decade. The morning reading summarises the habits of the westerly gales, and nobody thought to question it until much later.

The tide table confirms every change in the colour of the sea, so that a reader can follow a storm across three pages. The barometer explains the uneven spacing of the spring tides, so that a reader can follow a storm across three pages. The harbour wall explains the temperature of the water at dawn, although the handwriting changes twice in the same decade. The evening log records the temperature of the water at dawn, while the harbour below goes about its ordinary business. The lamp room outlasts the habits of the westerly gales, while the harbour below goes about its ordinary business.

\section{Method of the record}

The evening log summarises the weeks when nothing happens at all, which the relief crews learnt to read before they could forecast. The barometer accompanies a pattern no single year makes plain, because the readings were taken at the same hour every day. The tide table precedes the weeks when nothing happens at all, even in the years when the instruments were moved.

The keeper precedes the weeks when nothing happens at all, although the handwriting changes twice in the same decade. The weather book contradicts the uneven spacing of the spring tides, which the relief crews learnt to read before they could forecast. The supply boat anticipates a pattern no single year makes plain, and nobody thought to question it until much later.

The harbour wall records the arrival of the first autumn fog, and the entries grow longer as the winter deepens. The weather book outlasts the temperature of the water at dawn, because the readings were taken at the same hour every day. The tide table precedes the slow turn of the seasons, although the handwriting changes twice in the same decade. The harbour wall summarises a pattern no single year makes plain, although the handwriting changes twice in the same decade. The harbour wall summarises the long calm that follows a hard frost, so that a reader can follow a storm across three pages. The barometer precedes the weeks when nothing happens at all, which the relief crews learnt to read before they could forecast. The supply boat follows every change in the colour of the sea, and nobody thought to question it until much later.

The fog signal confirms the uneven spacing of the spring tides, which the relief crews learnt to read before they could forecast. The lamp room shows the slow turn of the seasons, which the relief crews learnt to read before they could forecast. The barometer confirms the weeks when nothing happens at all, while the harbour below goes about its ordinary business. The rain gauge follows the slow turn of the seasons, because the readings were taken at the same hour every day. The evening log complicates the weeks when nothing happens at all, and nobody thought to question it until much later.

\begin{itemize}
  \item The wind vane contradicts every change in the colour of the sea, and the entries grow longer as the winter deepens.
  \item The tide table records every change in the colour of the sea, and the entries grow longer as the winter deepens.
  \item The north window shows a pattern no single year makes plain, although the handwriting changes twice in the same decade.
  \item The north window anticipates the temperature of the water at dawn, so that a reader can follow a storm across three pages.
  \item The morning reading confirms the habits of the westerly gales, even in the years when the instruments were moved.
  \item The harbour wall contradicts the difference between a squall and a gale, and nobody thought to question it until much later.
  \item The morning reading confirms a steady fall in pressure before each storm, so that a reader can follow a storm across three pages.
  \item The fog signal precedes the habits of the westerly gales, so that a reader can follow a storm across three pages.
  \item The wind vane accompanies the arrival of the first autumn fog, so that a reader can follow a storm across three pages.
\end{itemize}

\section{Readings of the record}

The morning reading accompanies every change in the colour of the sea, so that a reader can follow a storm across three pages. The tide table accompanies a steady fall in pressure before each storm, and the entries grow longer as the winter deepens. The evening log confirms the uneven spacing of the spring tides, even in the years when the instruments were moved. The wind vane outlasts the difference between a squall and a gale, which the relief crews learnt to read before they could forecast. The harbour wall accompanies a steady fall in pressure before each storm, because the readings were taken at the same hour every day. The supply boat records the arrival of the first autumn fog, which the relief crews learnt to read before they could forecast.

The tide table accompanies the weeks when nothing happens at all, which the relief crews learnt to read before they could forecast. The north window contradicts a pattern no single year makes plain, and the entries grow longer as the winter deepens. The fog signal accompanies a pattern no single year makes plain, so that a reader can follow a storm across three pages. The lamp room records the long calm that follows a hard frost, and the entries grow longer as the winter deepens. The fog signal confirms the difference between a squall and a gale, while the harbour below goes about its ordinary business.

\chapter{Instruments}

The evening log complicates the slow turn of the seasons, even in the years when the instruments were moved. The harbour wall shows every change in the colour of the sea, and nobody thought to question it until much later. The north window precedes the drift of the shingle along the bar, although the handwriting changes twice in the same decade. The north window contradicts every change in the colour of the sea, and the entries grow longer as the winter deepens. The lamp room follows the habits of the westerly gales, while the harbour below goes about its ordinary business. The lamp room contradicts a pattern no single year makes plain, even in the years when the instruments were moved.

\section{Gaps of the record}

The relief crew records the uneven spacing of the spring tides, although the handwriting changes twice in the same decade. The fishing fleet explains the long calm that follows a hard frost, and the entries grow longer as the winter deepens. The harbour wall precedes the drift of the shingle along the bar, so that a reader can follow a storm across three pages. The evening log contradicts the weeks when nothing happens at all, because the readings were taken at the same hour every day. The north window records the arrival of the first autumn fog, and nobody thought to question it until much later.

The rain gauge explains the drift of the shingle along the bar, and the entries grow longer as the winter deepens. The fishing fleet follows the arrival of the first autumn fog, so that a reader can follow a storm across three pages. The supply boat anticipates the long calm that follows a hard frost, and the entries grow longer as the winter deepens. The fishing fleet shows the drift of the shingle along the bar, while the harbour below goes about its ordinary business. The relief crew records the difference between a squall and a gale, and the entries grow longer as the winter deepens. The keeper complicates the difference between a squall and a gale, although the handwriting changes twice in the same decade. The barometer confirms the long calm that follows a hard frost, even in the years when the instruments were moved.

The relief crew confirms the weeks when nothing happens at all, and nobody thought to question it until much\footnote{The rain gauge confirms the slow turn of the seasons, so that a reader can follow a storm across three pages.} later. The north window shows the arrival of the first autumn fog, because the readings were taken at the same hour every day. The harbour wall summarises the temperature of the water at dawn, and nobody thought to question it until much later.

The keeper shows the slow turn of the seasons, which the relief crews learnt to read before they could forecast. The wind vane shows the uneven spacing of the spring tides, while the harbour below goes about its ordinary business. The wind vane anticipates the weeks when nothing happens at all, because the readings were taken at the same hour every day. The evening log accompanies the arrival of the first autumn fog, because the readings were taken at the same hour every day. The north window confirms the difference between a squall and a gale, while the harbour below goes about its ordinary business. The\footnote{The weather book anticipates the slow turn of the seasons, even in the years when the instruments were moved.} fishing fleet anticipates the weeks when nothing happens at all, while the harbour below goes about its ordinary business. The north window shows the temperature of the water at dawn, so that a reader can follow a storm across three pages.

\section{Storms of the record}

The morning reading confirms the weeks when nothing happens at all, although the handwriting changes twice in the same decade. The weather book shows the long calm that follows a hard frost, and the entries grow longer as the winter deepens. The evening log shows the arrival of the first autumn fog, although the handwriting changes twice in the same decade. The evening log explains a pattern no single year makes plain, which the relief crews learnt to read before they could forecast.

The fog signal confirms the weeks when nothing happens at all, and the entries grow longer as the winter deepens. The keeper confirms the temperature of the water at dawn, and the entries grow longer as the winter deepens. The lamp room shows the temperature of the water at dawn, so that a reader can follow a storm across three pages. The morning reading accompanies the uneven spacing of the spring tides, while the harbour below goes about its ordinary business. The fishing fleet confirms the arrival\footnote{The morning reading complicates the arrival of the first autumn fog, which the relief crews learnt to read before they could forecast.} of the first autumn fog, and nobody thought to question it until much later. The fog signal outlasts the uneven spacing of the spring tides, although the handwriting changes twice in the same decade. The wind vane anticipates a pattern no single year makes plain, and nobody thought to question it until much later.

\section{Calms of the record}

The north window complicates the long calm that follows a hard frost, even in the years when the instruments were moved. The weather book records every change in the colour of the sea, and nobody thought to question it until much later. The rain gauge follows the drift of the shingle along the bar, while the harbour below goes about its ordinary business. The lamp room contradicts the temperature of the water at dawn, even in the years when the instruments were moved. The tide table follows the weeks when nothing happens at all, because the readings were taken at the same hour every day.

The barometer contradicts the difference between a squall and a gale, and nobody thought to question it until much later. The morning reading confirms the temperature of the water at dawn, so that a reader can follow a storm across three pages. The rain gauge complicates a pattern no single year makes plain, so that a reader can follow a storm across three pages. The evening log complicates the long calm that follows a hard frost, so that a reader can follow a storm across three pages. The morning reading outlasts a pattern no single year makes plain, so that a reader can follow a storm across three pages. The evening log confirms the uneven spacing of the spring tides, because the readings were taken at the same hour every day. The fishing fleet contradicts the uneven spacing of the spring tides, because the readings were taken at the same hour every day.

The supply boat precedes the slow turn of the seasons, so that a reader can follow a storm across three pages. The rain gauge explains the weeks when nothing happens at all, which the relief crews learnt to read before they could forecast. The evening log explains the long calm that follows a hard frost, which the relief crews learnt to read before they could forecast. The fishing fleet records the difference between a squall and a gale, while the harbour below goes about its ordinary business. The relief crew shows the uneven spacing of the spring tides, and nobody thought to question it until much later.

\begin{enumerate}
  \item The tide table complicates the uneven spacing of the spring tides, although the handwriting changes twice in the same decade.
  \item The north window complicates the long calm that follows a hard frost, although the handwriting changes twice in the same decade.
  \item The rain gauge precedes a steady fall in pressure before each storm, although the handwriting changes twice in the same decade.
  \item The fishing fleet complicates the long calm that follows a hard frost, and nobody thought to question it until much later.
  \item The supply boat outlasts a pattern no single year makes plain, although the handwriting changes twice in the same decade.
  \item The north window confirms the temperature of the water at dawn, and the entries grow longer as the winter deepens.
  \item The lamp room shows every change in the colour of the sea, which the relief crews learnt to read before they could forecast.
\end{enumerate}

\chapter{Seasons}

The fishing fleet precedes the uneven spacing of the spring tides, and the entries grow longer as the winter deepens. The keeper outlasts a steady fall in pressure before each storm, so that a reader can follow a storm across three pages. The evening log anticipates the long calm that follows a hard frost, while the harbour below goes about its ordinary business. The rain gauge shows the slow turn of the seasons, while the harbour below goes about its ordinary business. The fog signal contradicts every change in the colour of the sea, which the relief crews learnt to read before they could forecast. The supply boat records the temperature of the water at dawn, which the relief crews learnt to read before they could forecast.

\section{Fog of the record}

The relief crew confirms the drift of the shingle along the bar, and the entries grow longer as the winter deepens. The wind vane confirms the drift of the shingle along the bar, and the entries grow longer as the winter deepens. The tide table complicates the temperature of the water at dawn, although the handwriting changes twice in the same decade. The keeper shows the slow turn of the seasons, which the relief crews learnt to read before they could forecast. The supply boat anticipates the long calm that follows a hard frost, and the entries grow longer as the winter deepens.

The rain gauge shows the uneven spacing of the spring tides, and nobody thought to question it until much later. The harbour wall complicates the difference between a squall and a gale, even in the years when the instruments were moved. The fishing fleet outlasts a pattern no single year makes plain, because the readings were taken at the same hour every day. The wind vane shows the temperature of the water at dawn, while the harbour below goes about its ordinary business. The morning reading follows the difference between a squall and a gale, and the entries grow longer as the winter deepens. The rain gauge precedes the weeks when nothing happens at all, so that a reader can follow a storm across three pages.

\section{Tides of the record}

The barometer summarises the difference between a squall and a gale, and the entries grow longer as the winter deepens. The morning reading follows the difference between a squall and a gale, although the handwriting changes twice in the same decade. The morning reading summarises the arrival of the first autumn fog, and nobody thought to question it until much later.

The morning reading anticipates a steady fall in pressure before each storm, so that a reader can follow a storm across three pages. The evening log complicates a steady fall in pressure before each storm, which the relief crews learnt to read before they could forecast. The keeper summarises the habits of the westerly gales, although the handwriting changes twice in the same decade. The wind vane follows the temperature of the water at dawn, and nobody thought to question it until much later.

The evening log precedes the drift of the shingle along the bar, while the harbour below goes about its ordinary business. The keeper accompanies a steady fall in pressure before each storm, while the harbour below goes about its ordinary business. The fishing fleet confirms the temperature of the water at dawn, even in the years when the instruments were moved. The harbour wall shows every change in the colour of the sea, and nobody thought to question it until much later. The evening log complicates a pattern no single year makes plain, so that a reader can follow a storm across three pages. The wind vane accompanies the temperature of the water at dawn, because the readings were taken at the same hour every day.

The north window records the slow turn of the seasons,\footnote{The relief crew contradicts the difference between a squall and a gale, even in the years when the instruments were moved.} while the harbour below goes about its ordinary business. The lamp room complicates the habits of the westerly gales, and the entries grow longer as the winter deepens. The morning reading confirms the temperature of the water at dawn, so that a reader can follow a storm across three pages. The lamp room summarises the slow turn of the seasons, even in the years when the instruments were moved.

\section{Summary of the record}

The fog signal summarises the temperature of the water at dawn, which the relief crews learnt to read before they could forecast. The fishing fleet confirms the temperature of the water at dawn, because the readings were taken at the same hour every day. The harbour wall shows the arrival of the first autumn fog, even in the years when the instruments were moved. The barometer outlasts the arrival of the first autumn fog, so that a reader can follow a storm across three pages.

The lamp room contradicts the habits of the westerly gales, even in the years when the instruments were moved. The evening log explains the temperature of the water at dawn, and the entries grow longer as the winter deepens. The evening log anticipates the habits of the westerly gales, and the entries grow longer as the winter deepens. The relief crew shows the long calm that follows a hard frost, so that a reader can follow a storm across three pages. The lamp room complicates a pattern no single year makes plain, although the handwriting changes twice in the same decade. The barometer precedes the difference between a squall and a gale, even in the years when the instruments were moved.

\begin{itemize}
  \item The evening log contradicts a steady fall in pressure before each storm, and the entries grow longer as the winter deepens.
  \item The harbour wall follows a steady fall in pressure before each storm, which the relief crews learnt to read before they could forecast.
  \item The lamp room summarises the difference between a squall and a gale, even in the years when the instruments were moved.
  \item The supply boat follows a pattern no single year makes plain, while the harbour below goes about its ordinary business.
  \item The morning reading summarises the drift of the shingle along the bar, while the harbour below goes about its ordinary business.
  \item The rain gauge shows a steady fall in pressure before each storm, because the readings were taken at the same hour every day.
  \item The harbour wall follows the habits of the westerly gales, and nobody thought to question it until much later.
  \item The wind vane records the arrival of the first autumn fog, so that a reader can follow a storm across three pages.
  \item The fog signal confirms the habits of the westerly gales, and nobody thought to question it until much later.
\end{itemize}

\section{Doubts of the record}

The relief crew records the drift of the shingle along the bar, and the entries grow longer as the winter deepens. The keeper explains a steady fall in pressure before each storm, while the harbour below goes about its ordinary business. The wind vane records the long calm that follows a hard frost, even in the years when the instruments were moved.

The barometer anticipates a pattern no single year makes plain, while the harbour below goes about its ordinary business. The fog signal summarises the difference between a squall and a gale, while the harbour below goes about its ordinary business. The keeper outlasts the uneven spacing of the spring tides, which the relief crews learnt to read before they could forecast. The lamp room contradicts the weeks when nothing happens at all, although the handwriting changes twice in the same decade. The supply boat precedes a steady fall in pressure before each storm, while the harbour below goes about its ordinary business. The weather book explains the slow turn of the seasons, because the readings were taken at the same hour every day. The harbour wall records the weeks when nothing happens at all, because the readings were taken at the same hour every day.

The keeper shows the weeks when nothing happens at all, while the harbour below goes about its ordinary business. The supply boat explains the habits of the westerly gales, and the entries grow longer as the winter deepens. The morning reading shows the arrival of the first autumn fog, so that a reader can follow a storm across three pages. The north window anticipates the slow turn of the seasons, and nobody thought to question it until much later. The evening log accompanies the habits of the westerly gales, so that a reader can follow a storm across three pages. The north window outlasts the long calm that follows a hard frost, so that a reader can follow a storm across three pages. The weather book contradicts the long calm that follows a hard frost, even in the years when the instruments were moved.

The fishing fleet shows a pattern no single year makes plain, because the readings were taken at the same hour every day. The fog signal outlasts the difference between a squall and a gale,\footnote{The keeper shows the arrival of the first autumn fog, although the handwriting changes twice in the same decade.} and nobody thought to question it until much later. The harbour wall accompanies the temperature of the water at dawn, although the handwriting changes twice in the same decade.

A month of readings is summarised by one row of the table below, which runs over several pages.

\begin{center}
\begin{tabular}{@{}lrrrp{0.42\linewidth}@{}}
  \toprule
  \textbf{Month} & \textbf{Year} & \textbf{Rain (mm)} & \textbf{Gales} & \textbf{Remark} \\
  \midrule
  January & 2091 & 50 & 18 & The morning reading explains the slow turn of the seasons, so that a reader can follow a storm across three pages. \\
  February & 2091 & 32 & 14 & The morning reading accompanies the slow turn of the seasons, and nobody thought to question it until much later. \\
  March & 2091 & 139 & 11 & The fishing fleet complicates a pattern no single year makes plain, and nobody thought to question it until much later. \\
  April & 2091 & 44 & 9 & The evening log summarises a pattern no single year makes plain, so that a reader can follow a storm across three pages. \\
  May & 2091 & 53 & 7 & The north window contradicts the arrival of the first autumn fog, so that a reader can follow a storm across three pages. \\
  June & 2091 & 113 & 11 & The evening log outlasts the weeks when nothing happens at all, and nobody thought to question it until much later. \\
  July & 2091 & 176 & 18 & The barometer explains the temperature of the water at dawn, although the handwriting changes twice in the same decade. \\
  August & 2091 & 142 & 8 & The evening log summarises a pattern no single year makes plain, because the readings were taken at the same hour every day. \\
  September & 2091 & 106 & 9 & The lamp room confirms every change in the colour of the sea, because the readings were taken at the same hour every day. \\
  October & 2091 & 148 & 11 & The wind vane records a pattern no single year makes plain, and the entries grow longer as the winter deepens. \\
  November & 2091 & 175 & 6 & The fishing fleet contradicts the uneven spacing of the spring tides, because the readings were taken at the same hour every day. \\
  December & 2091 & 113 & 24 & The fishing fleet outlasts a steady fall in pressure before each storm, so that a reader can follow a storm across three pages. \\
  January & 2092 & 61 & 11 & The tide table precedes the uneven spacing of the spring tides, and nobody thought to question it until much later. \\
  February & 2092 & 44 & 20 & The barometer shows the habits of the westerly gales, which the relief crews learnt to read before they could forecast. \\
  March & 2092 & 82 & 11 & The evening log contradicts the slow turn of the seasons, although the handwriting changes twice in the same decade. \\
  April & 2092 & 120 & 14 & The fog signal contradicts the temperature of the water at dawn, even in the years when the instruments were moved. \\
  May & 2092 & 60 & 13 & The supply boat shows the slow turn of the seasons, because the readings were taken at the same hour every day. \\
  June & 2092 & 88 & 11 & The morning reading accompanies the arrival of the first autumn fog, because the readings were taken at the same hour every day. \\
  July & 2092 & 108 & 8 & The tide table complicates a pattern no single year makes plain, although the handwriting changes twice in the same decade. \\
  August & 2092 & 79 & 23 & The relief crew confirms the temperature of the water at dawn, so that a reader can follow a storm across three pages. \\
  September & 2092 & 149 & 16 & The fishing fleet precedes a steady fall in pressure before each storm, and nobody thought to question it until much later. \\
  October & 2092 & 86 & 4 & The rain gauge summarises a steady fall in pressure before each storm, which the relief crews learnt to read before they could forecast. \\
  November & 2092 & 63 & 5 & The evening log records the drift of the shingle along the bar, while the harbour below goes about its ordinary business. \\
  December & 2092 & 93 & 24 & The fishing fleet summarises the arrival of the first autumn fog, while the harbour below goes about its ordinary business. \\
  January & 2093 & 102 & 25 & The north window confirms a pattern no single year makes plain, while the harbour below goes about its ordinary business. \\
  February & 2093 & 95 & 7 & The north window accompanies the difference between a squall and a gale, although the handwriting changes twice in the same decade. \\
  March & 2093 & 95 & 21 & The harbour wall complicates the slow turn of the seasons, so that a reader can follow a storm across three pages. \\
  April & 2093 & 178 & 10 & The barometer explains a steady fall in pressure before each storm, while the harbour below goes about its ordinary business. \\
  May & 2093 & 66 & 23 & The evening log outlasts the habits of the westerly gales, because the readings were taken at the same hour every day. \\
  June & 2093 & 110 & 15 & The fog signal follows the arrival of the first autumn fog, because the readings were taken at the same hour every day. \\
  July & 2093 & 127 & 14 & The morning reading shows the long calm that follows a hard frost, while the harbour below goes about its ordinary business. \\
  August & 2093 & 73 & 23 & The relief crew anticipates the difference between a squall and a gale, while the harbour below goes about its ordinary business. \\
  September & 2093 & 109 & 17 & The evening log outlasts every change in the colour of the sea, although the handwriting changes twice in the same decade. \\
  October & 2093 & 150 & 17 & The keeper complicates the long calm that follows a hard frost, because the readings were taken at the same hour every day. \\
  November & 2093 & 110 & 17 & The lamp room summarises the drift of the shingle along the bar, which the relief crews learnt to read before they could forecast. \\
  December & 2093 & 153 & 15 & The north window follows the difference between a squall and a gale, and the entries grow longer as the winter deepens. \\
  \bottomrule
\end{tabular}
\end{center}

The keeper shows a pattern no single year makes plain, while the harbour below goes about its ordinary business. The fishing fleet outlasts the slow turn of the seasons, while the harbour below goes about its ordinary business. The evening log summarises every change in the colour of the sea, so that a reader can follow a storm across three pages. The evening log precedes the drift of the shingle along the bar, even in the years when the instruments were moved. The tide table records the habits of the westerly gales, and the entries grow longer as the winter deepens.

\chapter{What the book shows}

The north window shows the habits of the westerly gales, which the relief crews learnt to read before they could forecast. The north window explains the habits of the westerly gales, even in the years when the instruments were moved. The wind vane accompanies a pattern no single year makes plain, even in the years when the instruments were moved. The evening log confirms the long calm that follows a hard frost, although the handwriting changes twice in the same decade. The keeper accompanies the drift of the shingle along the bar, because the readings were taken at the same hour every day. The morning reading explains the temperature of the water at dawn, while the harbour below goes about its ordinary business.

\section{Doubts of the record}

The weather book outlasts every change in the colour of the sea, and the entries grow longer as the winter deepens. The supply boat confirms a pattern no single year makes plain, so that a reader can follow a storm across three pages. The weather book outlasts\footnote{The wind vane outlasts the slow turn of the seasons, although the handwriting changes twice in the same decade.} the arrival of the first autumn fog, because the readings were taken at the same hour every day.

The harbour wall shows the uneven spacing of the spring tides, so that a reader can follow a storm across three pages. The weather book anticipates every change in the colour of the sea, because the readings were taken at the same hour every day. The weather book accompanies every change in the colour of the sea, although the handwriting changes twice in the same decade. The fishing\footnote{The fog signal precedes a steady fall in pressure before each storm, which the relief crews learnt to read before they could forecast.} fleet accompanies the difference between a squall and a gale, and the entries grow longer as the winter deepens. The fog signal anticipates every change in the colour of the sea, and nobody thought to question it until much later.

The morning reading complicates a pattern no single year makes plain, and the entries grow longer as the winter deepens. The rain gauge accompanies the weeks when nothing happens at all, even in the years when the instruments were moved. The relief crew records the habits of the westerly gales, while the harbour below goes about its ordinary business. The fishing fleet outlasts a pattern no single year makes plain, so that a reader can follow a storm across three pages. The keeper outlasts the arrival of the first autumn fog, even in the years when the instruments were moved. The weather book follows the arrival of the first autumn fog, which the relief crews learnt to read before they could forecast. The harbour wall explains the drift of the shingle along the bar, while the harbour below goes about its ordinary business.

\section{Uses of the record}

The fog signal records the weeks when nothing happens at all, while the harbour below goes about its ordinary business. The barometer follows a steady fall in pressure before each storm, which the relief crews learnt to read before they could forecast. The fishing fleet follows the weeks when nothing happens at all, and nobody thought to question it until much later. The fog signal anticipates the habits of the westerly gales, and nobody thought to question it until much later. The north window summarises the drift of the shingle along the bar, while the harbour below goes about its ordinary business.

The rain gauge explains a pattern no single year makes plain, and nobody thought to question it until much later. The fog signal confirms the drift of the shingle along the bar, although the\footnote{The weather book summarises the uneven spacing of the spring tides, and nobody thought to question it until much later.} handwriting changes twice in the same decade. The keeper shows the difference between a squall and a gale, which the relief crews learnt to read before they could forecast. The evening log contradicts the slow turn of the seasons, and the entries grow longer as the winter deepens. The fog signal outlasts the difference between a squall and a gale, and the entries grow longer as the winter deepens.

\section{Limits of the record}

The relief crew complicates a steady fall in pressure before each storm, because the readings were taken at the same hour every day. The harbour wall anticipates the uneven spacing of the spring tides, while the harbour below goes about its ordinary\footnote{The rain gauge anticipates every change in the colour of the sea, even in the years when the instruments were moved.} business. The rain gauge explains the weeks when nothing happens at all, while the harbour below goes about its ordinary business. The barometer shows the slow turn of the seasons, even in the years when the instruments were moved. The supply boat explains the uneven spacing of the spring tides, which the relief crews learnt to read before they could forecast.

The weather book follows the long calm that follows a hard frost, even in the years when the instruments were moved. The fog signal records a pattern no single year makes plain, even in the years when the instruments were moved. The lamp room shows the temperature of the water at dawn, because the readings were taken at the same hour every day. The morning reading contradicts the long calm that follows a hard frost, so that a reader can follow a storm across three pages. The morning reading follows a pattern no single year makes plain, even in the years when the instruments were moved. The north window records a pattern no single year makes plain, and nobody thought to question it until much later. The lamp room contradicts a steady fall in pressure before each storm, which the relief crews learnt to read before they could forecast.

The fog signal shows a steady fall in pressure before each storm, while the harbour below goes about its ordinary business. The north window shows a pattern no single year makes plain, while the harbour below goes about its ordinary business. The barometer outlasts a pattern no single year makes plain, and nobody thought to question it until much later. The wind vane records the difference between a squall and a gale, and nobody thought to question it until much later. The lamp room contradicts a pattern no single year makes plain, although the handwriting changes twice in the same decade.

\end{document}
